\documentclass[../../../uem_mei_q.tex]{subfiles}

\begin{document}
    次の空欄\:\myboxb{　イ　}\,，\,\myboxb{　ウ　}\,，\,\myboxb{　エ　}\,，\,\myboxb{　オ　}\,，\\
    \myboxb{　カ　}\:に当てはまるものを指定された解答群の中から選び，%
    解答用紙の所定の欄の番号をマークせよ．なお，解答群から同じものを2回以上選んでもよい．%
    それ以外の空欄には，当てはまる0から9までの数字を解答用紙の所定の欄にマークせよ．%

    \vspace{.5\baselineskip}
    座標平面上で，方程式 $x^2-\myfracc{y^2}{4}=1$ によって定まる双曲線を$C$とする．%
    傾きが負となる$C$の漸近線を$l$とおくと，$l$の方程式は\\
    \hspace{3\zw}%
    $\displaystyle%
        y=-\:\myboxb{　ア　}\:x%
    $\\
    である．

    第1象限にある$C$上の点$\mathrm{P}(s,\;t)$をとると，%
    $t$は$s$の関数として $t=\:\myboxb{　イ　}\:s$ と表される．%
    点Pにおける$C$の接線を$m$とおくと，$m$の方程式は\\
    \hspace{3\zw}%
    $\displaystyle%
        \myboxb{　ウ　}\:x-\myfracc{\myboxb{　エ　}}{2}y=1%
    $\\        
    である．

    $x$軸と$m$との交点をQとおくと，Qの座標は%
    $\raisebox{-3pt}{$\Biggl($}\myfraca{1}{\myboxb{　オ　}},\;0\raisebox{-3pt}{$\Biggr)$}$%
    である．また，$l$と$m$の交点をRとおくと，Rの座標は \\
    $\left(\myfraca{1}{\myboxb{　カ　}},\;-\myfracc{\myboxa{　ア　}}{\myboxa{　カ　}}\right)$である．

    原点をOとし，三角形OQRの面積を$f(s)$とする．また，$x$軸上に点$\mathrm{S}(s,\;0)$をとり，%
    三角形PQSの面積を$g(s)$とする．\\
    このとき，\\
    \hspace{3\zw}%
    $\displaystyle%
        \lim_{s\to\infty}f(s)=\:\myboxb{　キ　}\:,\quad%
        \lim_{s\to\infty}f(s)g(s)=\myfracc{\myboxb{　ク　}}{\myboxb{　ケ　}}%
    $\\
    となる．

    \vspace{.5\baselineskip}
    \textsf{イ，ウ，エ，オ，カの解答群}

    {\renewcommand{\arraystretch}{1.3}%
        \vspace{-.1\baselineskip}%
        \hspace{\zw}%
        \begin{tabularx}{\dimexpr\linewidth-2\zw}{@{}LLL@{}}
            ⓪\; $2\sqrt{s^2-1}$ & ①\; $2s$ & ②\; $\sqrt{s^2-1}$ \\ [2pt]
            ③\; $s+2\sqrt{s^2-1}$ & ④\; $\sqrt{1-s^2}+s$ & ⑤\; $2\sqrt{1-s^2}$ \\ [2pt]
            ⑥\; $s$ & ⑦\; $\sqrt{1-s^2}$ & ⑧\; $s+\sqrt{s^2-1}$ \\ [2pt]
            ⑨\; $\sqrt{1-s^2}-s$
        \end{tabularx}%
    }
\end{document}
